\documentclass{ximera}

\title{Derivative Rules (So Far) Activity Facilitator Guide}
\author{MATH 425: Calculus I}

\begin{document}
\begin{abstract}
    Working with the peers in your group, solve the following problems. Make sure to show and justify all your work. Make sure everyone in the group understands the solution and participates. Be prepared to report your answers to the whole class. 
    %A default abstract of an almost empty course.
\end{abstract}
\maketitle

\textbf{Student Learning Outcomes}\\
Students will be able to: 
    \begin{itemize}
        \item Take derivatives using the power, constant multiple, sum, and difference rules for functions involving these operations.
        \item Find the tangent line of a function involving power, constant multiple, sum, and difference rules.
        \item Simplify derivative definitions involving combinations of functions in order to justify derivative rules.
    \end{itemize}

\textbf{Prerequisite Knowledge:}  This activity assumes that students have been introduced to the derivative rules, excepting the product, quotient and chain rules.  Algebraic simplificaiton may be needed in order to take derivatives of some parts based on the given rules.  \\

The rules were purposefully not given on this activity.  This is a great opportunity to have students pull out their notes in order to have the rules in front of them as they work through this activity.  If they do not have the rules written down, please encourage them to do so.  \\

This activity may work well as a ``choose your own adventure'': groups that are comfortable with the rules can start with exercises 3-4, then return to 1-2 for the practice.  Groups that are less confident with the rules can start with 1-2 for reinforcement, then work on 3-4 as time allows.  


\begin{exercise}
Calculate the derivative for the following functions using the derivative rules we've discussed in class so far: 
\begin{enumerate}
    \item $s(t)=\frac{1}{t}-\frac{3}{t^6}$\\
    \textcolor{blue}{$s'(t)=-t^{-2}+18t^{-7}$}

    \item $f(x)=\sqrt{3}x^2-\sqrt{3}$
    \textcolor{blue}{ $f'(x)=2\sqrt{3}x$}
    
    \item $g(w)=w^5-\frac{1}{w^\frac{3}{2}}$
    \textcolor{blue}{ $g'(w)=5w^4+\frac{3}{2}w^{-\frac{5}{2}}$}
    
    \item $f(z)=\frac{2\sqrt{z}}{3}$ 
    \textcolor{blue}{ $f'(x)= \frac{2}{3}\cdot \frac12 z^{-1/2}=\frac13 z^{-1/2}$}
    
    \item $g(x)=-7x^{0.25}$
    \textcolor{blue}{ $g'(x)=-\frac{7}{4}x^{-0.75}$}
    
    \item $v(t)=\pi^2$
    \textcolor{blue}{ $v'(t)=0$}
    
    \item $p(t)=\frac{2t^2-3t+1}{\sqrt{t}}$ \textcolor{blue}{ $=2t^{\frac32}-3t^{\frac12}+t^{-\frac12}$}\\
    \textcolor{blue}{ $p'(t)=3t^{\frac12}-\frac32 t^{-\frac12}-\frac12t^{-\frac32}$}\\
    \textcolor{orange}{Note that this problem can be done with the quotient rule, but the intended method at this poin in the course is to simplify first.}

    \item $s(t)=\frac{1}{t}-\frac{3}{t^6}+e^5$\\
    \textcolor{blue}{ $s'(t)=-t^{-2}+18t^-7+0$}
\end{enumerate}
\end{exercise}

\begin{exercise}
    Find the equation of the line tangent to the graphs of the following functions at the given point.  Use Desmos to graph the function and your tangent line and check your work.
        \begin{enumerate}
    %\item $f(x)=e^x$ at $(0,1)$\\
            \item $g(x)=x^6-24x$ at $(2, 16)$\\
            \textcolor{blue}{ $g'(x)=6x^5-24$ so $g'(2)=6(2)^5-24=192-24=168$\\
            So the tangent line is $y-16=168(x-2)$ or $y=168x-320$}
            \item $h(x)=\sqrt[3]{x}-\frac{3}{x}$ at $(1, -2)$\\
            \textcolor{blue}{ $h'(x)=\frac{1}{3}x^{-2/3}+3x^{-2}=\frac{1}{3\sqrt[3]{x^2}}+\frac{3}{x^2}$ so $h'(1)=\frac{1}{3\sqrt[3]{1^2}}+\frac{3}{1^2}=\frac{1}{3}+3=\frac{10}{3}$\\
            So the tangent line is $y+2=\frac{10}{3}(x-1)$ or $y=\frac{10}{3}x-\frac{16}{3}$}
        \end{enumerate}
        \textcolor{orange}{Tangent lines may be left in point-slope or point-intercept form.}
\end{exercise}

\textcolor{orange}{Problems 3 and 4 are intended to be justifications of derivative rules, rather than detailed proofs.  The scaffolding in problem 3 is intended to step students through, but can be a little wordy for some students.  If students are getting caught up in the wording, encourage them to play with the derivative rules and try to get the result on their own.  The difference rule can also be justified in a similar way if groups are interested.  The power rule is more complicated because of the needed expansion and/or factoring of the $(x+h)^n-x^n$ but advanced students with help may be able to justify it as well.}

\begin{exercise}
  The following problem will step you through proving the sum rule for derivatives. Our goal will be to show that $$\frac{d}{dx}[(f+g)(x)]=\frac{d}{dx}[f(x)]+\frac{d}{dx}[g(x)]$$
    \begin{enumerate}
        \item Write $(f+g)(x+h)$ in terms of $f(x)$ and $g(x)$. \\
        \textcolor{blue}{ $(f+g)(x+h)=f(x+h)+g(x+h)$}
        \item Use your answer from (a) to write the derivative definition for $(f+g)(x)$.\\
        \textcolor{blue}{ $\displaystyle \lim_{h\rightarrow 0}\frac{(f+g)(x+h)-(f+g)(x)}{h}=\lim_{h\rightarrow 0}\frac{f(x+h)+g(x+h)-f(x)-g(x)}{h}$}
        \item Group your numerator of the derivative definition into $f$ terms and $g$ terms.\\
        \textcolor{blue}{$\displaystyle =\lim_{h\rightarrow 0}\frac{(f(x+h)-f(x))+(g(x+h)-g(x))}{h}$}
        \item Use the opposite of finding a common denominator to split your fraction into two fractions (hint: both should look like the difference quotient $\frac{F(x+h)-F(x)}{h}$).\\
        \textcolor{blue}{ $\displaystyle =\lim_{h\rightarrow 0} \frac{f(x+h)-f(x)}{h}+\frac{g(x+h)-g(x)}{h}$}
        \item Use the sum property for limits to separate your limit into the sum of two limits.\\
        \textcolor{blue}{ $\displaystyle =\lim_{h\rightarrow 0} \frac{f(x+h)-f(x)}{h}+\lim_{h\rightarrow 0}\frac{g(x+h)-g(x)}{h}$}
        \item Substitute a derivative notation for any derivative definitions  that show up. $\square$ ($\leftarrow$ this is a symbol traditionally used to show a proof is complete.) \\
        \textcolor{blue}{ $=\displaystyle f'(x)+g'(x)$}
    \end{enumerate}
\end{exercise}

\begin{exercise}
    Using a similar method to (4), show that $$ \frac{d}{dx}(cf(x))=c\frac{d}{dx}f(x)$$ by manipulating the definition of the derivative.\\
    \textcolor{blue}{ $\displaystyle \frac{d}{dx}(cf(x))=\lim_{h\rightarrow 0} \frac{cf(x+h)-cf(x)}{h}=\lim_{h\rightarrow 0}\frac{c(f(x+h)-f(x))}{h}$\\
    $\displaystyle =\lim_{h\rightarrow 0}c\frac{f(x+h)-f(x)}{h}=c\lim_{h\rightarrow 0}\frac{f(x+h)-f(x)}{h}=cf'(x)$}
\end{exercise}




\end{document}